\section{Heavy-ball Momentum Convergence}
\label{sec:momentum}
We now turn our attention to the analogous dynamics in an inertial setting that models heavy-ball momentum-based optimization.
Upon introducing momentum, the system no longer closes on the position covariance of the neuron weights \((u,v)\) alone: one must instead consider the phase space of positions and velocities tracked jointly.
We show that, in terms of the \emph{phase covariance}, dynamics are again closed for linear networks.
We derive conservation identities describing this phase geometry; energy dissipation and control of position deformation give the convergence certificate below.

\paragraph{Dynamics and phase closure}
Inertia adds velocity as a state variable, so the natural moment over which the dynamics close is the joint position-velocity covariance.
For $W:=\binom UV$, predictor $\Svu:=VU^\top/n$, and damping $\gamma>0$,
write $E:=\nabla\ell(\Svu)$. The factor dynamics are
\begin{equation}\label{eq:finite-momentum}
 \ddot W+\gamma\dot W+G(\Svu)W=0,\qquad
 G(\Svu):=\begin{psmallmatrix}0&E^\top\\E&0\end{psmallmatrix}.
\end{equation}
Write $(P,Q):=(\dot U,\dot V)$, so the columns of $W$ and $\dot W$ are
positions $x_i:=(u_i,v_i)$ and velocities $y_i:=(p_i,q_i)$ respectively.
In this section $\Phase:=n^{-1}\binom W{\dot W}\binom W{\dot W}^{\!\top}$
denotes the full $2m$-dimensional phase covariance. Its blocks are
$\Spos:=WW^\top/n$, $\Svel:=\dot W\dot W^\top/n$, and
$\Xi:=W\dot W^\top/n$; the predictor remains the $vu$ block of $\Spos$.
With $G_t:=G(\Svu_t)$, the closed equation is
\begin{equation}\label{eq:phase}
 \dot\Phase_t=\Mob_t\Phase_t+\Phase_t\Mob_t^\top,\qquad
 \Phase_t=\begin{psmallmatrix}\Spos_t&\Xi_t\\\Xi_t^\top&\Svel_t\end{psmallmatrix},\qquad
 \Mob_t:=\begin{psmallmatrix}0&\Id_m\\-G_t&-\gamma\Id_m\end{psmallmatrix}.
\end{equation}
Appendix~\ref{app:momentum-proof} gives the particle equations, measure
formulation, and blockwise expansion.

\paragraph{Conformal laws and mechanical energy}
Without damping, the fundamental matrix generated by $\Mob_t$ in \eqref{eq:phase} is symplectic;
damping rescales the associated bilinear form exponentially in time. 
\begin{proposition}[Conformal phase conservation laws]\label{prop:phase-laws}
Let $\Jcan:=\begin{psmallmatrix}0&\Id_m\\-\Id_m&0\end{psmallmatrix}$.
The phase fundamental matrix \(\dot\Psi_t=\Mob_t\Psi_t\),
\(\Psi_0=\Id_{2m}\) generated by $\Mob_t$ satisfies
\[
 (\Psi_t)^\top \Jcan\Psi_t=e^{-\gamma t}\Jcan,\qquad
 \Phase_t=\Psi_t\Phase_0(\Psi_t)^\top.
\]
Consequently, the spectrum of \(e^{\gamma t}\Jcan\Phase_t\)
and all its trace powers are constant. At the characteristic level,
these identities imply that, for any two trajectories of
the same flow, \(z_a:=(x_a,y_a)\), \(z_b:=(x_b,y_b)\),
\(e^{\gamma t}z_a^\top \Jcan z_b\) is conserved.
\end{proposition}

\begin{comment}
\emph{Proof sketch.}
Symmetry of $G_t$ gives $\Mob_t^\top\Jcan+\Jcan\Mob_t=-\gamma\Jcan$.
Differentiating $\Psi_t^\top\Jcan\Psi_t$ proves the conformal identity;
the identity $\Phase_t=\Psi_t\Phase_0\Psi_t^\top$ then shows that $e^{\gamma t}\Jcan\Phase_t$
is similar to $\Jcan\Phase_0$.\qed
\end{comment}

Special spectral values permit reduced phase dynamics, we describe these in Appendix~\ref{app:intrinsic}.
As before, write $\Delta_t:=\ell(\Svu_t)-\ell_\star$ for the predictor loss gap. Define $\mathcal{K}_t:=\tr(\Svel_t)/2$ and $\cE_t:=\Delta_t+\mathcal{K}_t$ as the kinetic and mechanical excess energies. Differentiating $\cE_t$ along \eqref{eq:phase} gives $\dot{\cE}_t=-2\gamma\mathcal{K}_t$.
This identity gives global existence of the solution to \eqref{eq:phase} from every $\Sigma_0 \succeq 0$ (Proposition~\ref{prop:global-phase}).
The remaining task is to derive spectral lower bounds for $\Spos_t$.

\paragraph{Global convergence}
We seek an initialization condition on $\Sigma_0$ that keeps $\Spos_t$ uniformly elliptic for the whole trajectory.
Let $a_0:=\lambda_{\min}(\Spos_0)>0$ and $b_0:=\lambda_{\max}(\Spos_0)$.
A deformation budget $s>0$ specifies a spectral interval $[a_s,b_s]:=[a_0e^{-2s},b_0e^{2s}]$ for $\Spos_t$.
For $\cE_0>0$, put $\rho_s:=4b_sL+2\sqrt{2L\cE_0}$ and define
\begin{equation}\label{eq:certificate}
 \lambda_s(\gamma):=\tfrac{a_s\kappa\gamma}{\rho_s+2\gamma^2},\qquad
 \mathfrak S_s(\gamma):=\tfrac4{\lambda_s(\gamma)}
             \sqrt{\tfrac{3\cE_0}{2a_s}}.
\end{equation}
Here, $\lambda_s$ is the exponential decay rate of $\cE_t$, and $\mathfrak S_s$ bounds the resulting total deformation in $\Spos_t$, conditional on the spectral interval $[a_s, b_s]$.

\begin{theorem}[Global linear rate from a phase-energy certificate]\label{thm:momentum}\label{thm:bootstrap-unconditional-rate}
Under Assumption~\ref{ass:loss}, let $\Phase_0\succeq0$ with
$\Spos_0\succ0$ and $\cE_0>0$. If
$\mathfrak S_s(\gamma)<s$ for some $s>0$, then for all $t\ge0$,
\[
 a_s\Id_m\preceq\Spos_t\preceq b_s\Id_m,\qquad
 \cE_t\le3\cE_0 e^{-\lambda_s(\gamma)t}.
\]
The covariance $\Spos_t$ and any associated factors
$(U_t,V_t)$ evolving under \eqref{eq:finite-momentum} converge exponentially with rate $\lambda_s/2$ to limits with optimal predictor and $\Spos_\infty\succ0$; $\Svel_t$ tends to zero.
\end{theorem}

\begin{comment}
\paragraph{Proof sketch}
\emph{Conditional energy decay.}
Global existence of $\Phase_t$ follows from the energy identity $\dot{\cE}_t=-2\gamma\mathcal{K}_t$.
When the eigenvalues of $\Spos_t$ lie in $[a_s,b_s]$, the PL inequality gives a lower bound on the mean squared particle force in terms
of $\Delta_t$. For sufficiently small $\epsilon>0$, we show that $\cE_t+\epsilon\dot\Delta_t$ is a Lyapunov function comparable to $\cE_t$, whose exponential decay yields the bound on $\cE_t$ (Lemma~\ref{lem:conditional}).

\emph{Controlled deformation.}
Whenever $\Spos_t\succ0$, define $D_t:=\Xi_t^\top(\Spos_t)^{-1}$.
Then $\dot\Spos_t=D_t\Spos_t+\Spos_tD_t^\top$ and
$\norm{D_t}_{\op}\le\sqrt{2\mathcal{K}_t/a_s}$ while the eigenvalues of
$\Spos_t$ remain in $[a_s,b_s]$. Using $\mathcal{K}_t\le\cE_t$ and the
energy bound gives $\int_0^t\norm{D_r}_{\op}\,\dd r\le\mathfrak S_s(\gamma)<s$.

\emph{Preserving the spectral bounds.}
Gr\"onwall's inequality and the integral bound above prevent the eigenvalues
of $\Spos_t$ from leaving $[a_s,b_s]$, so the energy bound and convergence
estimates hold globally. Appendix~\ref{app:momentum-proof} gives the full
proof and explicit convergence bounds; Appendix~\ref{app:momentum-details}
gives a sharper conditional energy bound.\qed
\end{comment}

\begin{corollary}[Two-scale white initialization]
\label{cor:two-scale-white-window}
Under Assumption~\ref{ass:loss}, let
\(\Spos_0=\sigma_{vu}^2\Id_m\), \(\Xi_0=0\), and
\(\Svel_0\succeq0\) have trace \(m\sigma_{qp}^2\), with
\(\sigma_{vu}>0\), \(\sigma_{qp}\ge0\).
Then \(\cE_0=\Delta+m\sigma_{qp}^2/2\), where
\(\Delta:=\ell(0)-\ell_\star\).
If \(\cE_0>0\), the certificate is
% \begin{equation}\label{eq:white-direct}
\[
 2\gamma^2-
 \tfrac{s\kappa\sigma_{vu}^3e^{-3s}}
 {4\sqrt{3\cE_0/2}}\gamma
 +4L\sigma_{vu}^2e^{2s}
 +2\sqrt{2L\cE_0}<0.
\]
% \end{equation}
For fixed \(s,\gamma>0,\Delta,\sigma_{qp},m,L,\kappa\), sufficiently
large \(\sigma_{vu}\) satisfies this condition.
\end{corollary}

Appendix~\ref{app:large-position-scale} gives explicit bounds on the admissible $\gamma$ interval and analyzes the large-$\sigma_{vu}$ asymptotics.

\paragraph{Numerical view of the bootstrap certificate}
\begin{wrapfigure}{r}{0.34\linewidth}
\centering
\resizebox{\linewidth}{!}{\subimport{figures/}{momentum_bootstrap_heatmap_notebook.pgf}}
\footnotesize
\captionsetup{font=footnotesize,skip=2pt}
\caption{Two-scale white momentum initialization. Curves show
$\mathfrak G_s(\gamma)=0$ for the color-bar scale $\sigma_{vu}$;
their interiors are certified damping windows.}
\label{fig:momentum-bootstrap-heatmap}\label{fig:momentum-window}
\end{wrapfigure}
Writing $\mathfrak G_s(\gamma):=\mathfrak S_s(\gamma)-s$, Figure~\ref{fig:momentum-bootstrap-heatmap} shows the boundaries
$\mathfrak G_s(\gamma)=0$ for $1\le\sigma_{vu}\le10^3$, with $(d_u,d_v)=(3,2)$,
$\kappa=L=1$, $\sigma_{qp}=1$, and $\ell(0)-\ell_\star=1$.
The regions where $\mathfrak G_s(\gamma)<0$ specify the pairs $(s,\gamma)$ for which Theorem~\ref{thm:momentum} guarantees spectral bounds and exponential energy decay.
For the two-scale initialization, sufficiently large $\sigma_{vu}$ yields a nonempty interval of $\gamma$ satisfying Corollary~\ref{cor:two-scale-white-window}.
For fixed initialization and damping, a smaller admissible $s$ gives a
tighter tube and faster rate; Figure~\ref{fig:momentum-bootstrap-lambda-multi-sigma}
illustrates the limiting rate at the feasibility boundary analyzed in
Appendix~\ref{app:large-position-scale}.

Figure~\ref{fig:momentum-empirical-vs-theory} compares empirical losses with the upper bounds on $\cE_t$ from Corollary~\ref{cor:explicit-tightened-conditional-momentum-rate}.
Each upper bound curve uses an explicit, approximately optimal rate certified by Theorem~\ref{thm:momentum} (see Proposition \ref{prop:optimal-tightened-conditional-momentum-rate}); no rate is fitted
to the empirical trajectory. All displayed losses stay below their upper bounds.
Notably in this inertial setting, the predictor loss may oscillate even while the mechanical energy decreases.

\begin{figure}[!htbp]
\centering
\resizebox{\linewidth}{!}{\subimport{figures/}{momentum_empirical_vs_tightened_theory_by_sigma.pgf}}
\captionsetup{font=footnotesize,skip=2pt}
\caption{
Empirical predictor-loss trajectories (solid) and the corresponding upper bounds
on mechanical excess energy (dotted) for two-scale white initializations with varying
\(\sigma_{vu}\) and damping \(\gamma\): \textbf{(a)} $\sigma_{vu}=200$,
\textbf{(b)} $\sigma_{vu}=500$, and \textbf{(c)} $\sigma_{vu}=1000$.
Each theoretical curve uses the explicit, approximately optimal rate on a tube certified by
Theorem~\ref{thm:bootstrap-unconditional-rate}.
}
\label{fig:momentum-empirical-vs-theory}\label{fig:momentum-envelopes}
\end{figure}
\WFclear

%%% Local Variables:
%%% mode: LaTeX
%%% TeX-master: "../main"
%%% End:
